% ======================================================================
\chapter{Scaling Laws in Practice}\label{ch:scalinglaws}

\Cref{sec:pretrain:scaling} gave us a planning rule: spend a compute budget $C$ on a model of
$N\approx\sqrt{C/120}$ parameters trained on $D\approx 20N$ tokens. That rule is the
\emph{output} of a scaling study somebody else ran, on their data, architecture, optimizer,
tokenizer and hardware. Change any of these and the constants move. This chapter is about
running the study yourself: how to turn a few dozen small training runs into a defensible
prediction for one large run you can only afford to do once.

For a systems engineer, a scaling study is an experiment-design problem under a hard resource
budget. The runs are jobs on a shared cluster. They have queueing delays, reservations, timeouts
and failures, and their wall-clock cost depends on how well each shape uses the hardware.
The analysis is numerical optimisation on a handful of noisy points, and the deliverable is an
extrapolation over several orders of magnitude, where small errors in an exponent become large
errors in parameter count. We will treat all three layers (accounting, fitting and
orchestration) with code, and finish with a complete budgeted study against a synthetic training
service that you can run on a laptop.

\begin{figure}[t]
\centering
\begin{tikzpicture}[
  stage/.style={draw, rounded corners=2pt, minimum width=2.35cm, minimum height=0.95cm,
    align=center, font=\footnotesize, fill=#1},
  stage/.default=fcompute!15,
  num/.style={circle, fill=sysblue, text=white, font=\scriptsize\bfseries, inner sep=1.2pt},
  arr/.style={-{Stealth[length=2mm]}, thick, sysblue!70}]
\node[stage] (a) at (0,0) {Why power\\laws?};
\node[stage] (b) at (2.8,0) {Compute\\accounting};
\node[stage] (c) at (5.6,0) {Kaplan vs\\Chinchilla};
\node[stage] (d) at (8.4,0) {Three fitting\\methods};
\node[stage=keygreen!15] (e) at (11.2,0) {LR, batch,\\$\mu$P};
\node[stage=keygreen!15] (f) at (11.2,-1.8) {WSD: cheap\\studies};
\node[stage=keygreen!15] (g) at (8.4,-1.8) {Data, inference,\\downstream};
\node[stage=napkin!20] (h) at (5.6,-1.8) {Budgeted study\\design};
\node[stage=napkin!20] (i) at (2.8,-1.8) {Uncertainty \&\\wall-clock};
\node[stage=napkin!20] (j) at (0,-1.8) {Running real\\sweeps};
\foreach \x/\n in {a/1,b/2,c/3,d/4,e/5,f/6,g/7,h/8,i/9,j/10}
  \node[num] at (\x.north west) {\n};
\draw[arr] (a) -- (b); \draw[arr] (b) -- (c); \draw[arr] (c) -- (d); \draw[arr] (d) -- (e);
\draw[arr] (e) -- (f); \draw[arr] (f) -- (g); \draw[arr] (g) -- (h); \draw[arr] (h) -- (i);
\draw[arr] (i) -- (j);
\end{tikzpicture}
\caption{Roadmap of this chapter. Steps 1--4 build the theory and the fitting machinery;
steps 5--7 cover the hyperparameter and data questions a real study must answer; steps 8--10 are
the systems side: planning, executing and trusting a study under a wall-clock budget.}
\label{fig:scalinglaws:roadmap}
\end{figure}

% ======================================================================
\section{Why do power laws appear?}\label{sec:scalinglaws:why}

A \emph{scaling law} is an empirical relation between a resource (data $D$, parameters $N$,
compute $C$) and a quality metric (usually held-out loss) that is a straight line on log-log axes
over a wide range:
\begin{equation}
L(x) \;=\; E + A\,x^{-\alpha}
\quad\Longleftrightarrow\quad
\log\bigl(L(x)-E\bigr) = \log A - \alpha\log x.
\label{eq:scalinglaws:pl}
\end{equation}
The constant $E$ is the irreducible loss: the entropy of the text itself, which no model can
beat. The exponent $\alpha$ says how fast extra resources pay off.

Why a power law rather than an exponential or a logarithm? The cleanest intuition comes from
statistics. Estimating the mean of $n$ samples from $\mathcal N(\mu,\sigma^2)$ has
mean-squared error $\sigma^2/n$, so $\log\text{error} = -\log n + 2\log\sigma$: a power law with
slope $-1$. Most parametric estimators share this $1/n$ rate. Flexible, nonparametric
estimators are slower. If you estimate an arbitrary smooth function on a $d$-dimensional
input by averaging over little boxes, you trade the number of boxes against the samples per box,
and the error falls only as $n^{-c/d}$. Language models behave like the second case: measured
data exponents are far smaller than 1 (Kaplan et al.\ report roughly $0.1$ for loss versus
data). One hypothesis is that the slope reflects the \emph{intrinsic dimension} of the data
manifold (Bahri et al., 2021), but this is an argument for why power laws are plausible, not a
derivation of the numbers. Treat the exponents as things to be measured.

\begin{figure}[t]
\centering
\begin{tikzpicture}
\begin{loglogaxis}[width=0.62\linewidth, height=5cm, xmin=10, xmax=1e6, ymin=1e-6, ymax=2,
  xlabel={samples $n$}, ylabel={estimation error}, label style={font=\small},
  tick label style={font=\scriptsize},
  legend style={font=\scriptsize, draw=none, at={(1.02,0.5)}, anchor=west}]
\addplot[very thick, sysblue, domain=10:1e6, samples=10] {1/x};
\addplot[very thick, napkin, domain=10:1e6, samples=10] {x^(-0.5)};
\addplot[very thick, labpurple, domain=10:1e6, samples=10] {x^(-0.1)};
\legend{parametric: $n^{-1}$, nonparametric $d{=}2$: $n^{-1/2}$, nonparametric $d{=}10$: $n^{-1/10}$}
\end{loglogaxis}
\end{tikzpicture}
\caption{Error rates that are polynomial in the sample size are straight lines on log-log axes.
The slope is the exponent; higher-dimensional problems give shallower lines. LM data
exponents look like the shallow end of this picture (analytic curves, constants set to 1).}
\label{fig:scalinglaws:rates}
\end{figure}

\begin{pitfall}[A straight line only after subtracting $E$]
With an irreducible term, $\log L$ is \emph{not} linear in $\log x$: the curve flattens as
$L\to E$. Fitting a line to $\log L$ directly underestimates $\alpha$ at the small-loss end and
overestimates the payoff of scale when you extrapolate into the region where $E$ dominates.
Either fit $E$ jointly (\cref{sec:scalinglaws:method3}) or stay in the regime where $L-E$ is
still large compared with the noise.
\end{pitfall}

Two more facts from the early literature shape how studies are run. First, \emph{distribution
shift} mostly changes the offset of a scaling curve, not its slope (Hashimoto, 2021): better data
moves the whole line down. Second, data scaling has been observed across modalities and
architectures (Hestness et al., 2017), so the method generalises even when the constants do not.

% ======================================================================
\section{What scaling laws are used for}\label{sec:scalinglaws:uses}

The scaling-law design procedure has three steps: train a family of small models, fit a law
to each candidate choice, and pick the choice whose law is best \emph{at the target scale}. The
last clause matters: a modification that helps at $10^{18}$ FLOPs can have a shallower slope and
lose at $10^{23}$. Kaplan et al.\ (2020) used this lens to compare transformers with LSTMs (the
LSTM curves bend away at large $N$), Adam with SGD, and different depth-to-width ratios. Three
lessons from that era are still the default assumptions of modern studies:
\begin{itemize}
  \item \textbf{Shape matters little within a broad basin.} At fixed non-embedding parameter
        count, loss is within a few percent over a wide range of aspect ratios $\dmodel/L$.
        Studies therefore fix the aspect ratio (or a rule for it) and scale one ``size'' knob.
  \item \textbf{Embedding parameters are different.} Loss tracks \emph{non-embedding}
        parameters much better than total parameters; at small scale the embedding table can be
        most of the model and would distort the fit (\cref{sec:scalinglaws:accounting}).
  \item \textbf{Hyperparameters that change with scale must be modelled.} Learning rate and
        batch size have scale-dependent optima (\cref{sec:scalinglaws:hparams}). Using one fixed
        value for every size silently bends the curve.
\end{itemize}

% ======================================================================
\section{Compute accounting: which $N$, which $C$?}\label{sec:scalinglaws:accounting}

\Cref{sec:pretrain:scaling} derived $C\approx 6ND$: two FLOPs per parameter per token in the
forward pass and four in the backward pass. A scaling study needs this more precisely, because
the fits are only as consistent as the accounting.

For a dense pre-norm transformer with SwiGLU MLPs, width $d$, $L$ layers, MLP width $\dff$,
vocabulary $V$ and context $T$, the matmul weights per layer are $4d^2$ (Q, K, V, O) plus
$3d\dff$ (gate, up, down). With the common $\dff = \tfrac83 d$ this gives the familiar
\begin{equation}
N_{\text{ne}} \;=\; L\,(4d^2 + 3d\dff) \;=\; 12\,L\,d^2 \qquad(\dff=\tfrac83d),
\label{eq:scalinglaws:12ld2}
\end{equation}
non-embedding parameters. Training FLOPs per token add the LM-head matmul ($Vd$ weights) and the
attention score and value products, which have no weights but cost $2LTd$ forward FLOPs per
token:
\begin{equation}
\frac{C}{D} \;\approx\; 6\,(N_{\text{ne}} + Vd) \;+\; 6\,L\,T\,d.
\label{eq:scalinglaws:fpt}
\end{equation}
The input embedding is a gather and costs no FLOPs. \Cref{lst:scalinglaws:runconfig} is the
accounting class every tool in this chapter uses.

\begin{lstlisting}[caption={Accounting for one run (from \code{llmsys/scaling.py}). Everything in this chapter---planning, fitting, the synthetic API---derives $N$ and $C$ from this one place.},label={lst:scalinglaws:runconfig}]
@dataclass(frozen=True)
class RunConfig:
    n_layer: int; d_model: int; n_head: int; d_ff: int; tokens: int
    batch_size: int = 128; lr: float = 3e-3; seq_len: int = 512
    vocab: int = 32_000; tie_embeddings: bool = False

    def non_embedding_params(self):
        d = self.d_model
        return self.n_layer * (4 * d * d + 3 * d * self.d_ff)

    def flops_per_token(self):
        matmul_params = self.non_embedding_params() + self.vocab * self.d_model
        return 6.0 * matmul_params + 6.0 * self.n_layer * self.seq_len * self.d_model

    def train_flops(self):
        return self.flops_per_token() * self.tokens
\end{lstlisting}

\begin{napkinmath}[How much of a small model is embedding?]
A model with $L=8$, $d=512$, $\dff=1408$ and $V=32{,}000$ has
$N_{\text{ne}} = 8(4\cdot512^2 + 3\cdot512\cdot1408) \approx 25.7$M, but its untied input
and output embeddings hold $2\cdot32{,}000\cdot512\approx 32.8$M. More than half the weights
are embeddings, and the LM head alone adds $16.4$M to the FLOP count, about $64\%$ of the
non-embedding work. If you fit a law with total parameters at this scale and use it to
extrapolate to a 1B model, where embeddings are a few percent of the weights, your $N$ axis
changes meaning along the way. Pick one definition (non-embedding is the usual choice), use it
for every run, and convert at the end.
\end{napkinmath}

\begin{systemslens}[A wall-clock budget is not a FLOP budget]
Clusters and training services charge in device-hours, not FLOPs. A run of $C$ FLOPs on a
device with peak $P$ takes
\[
t \;=\; \frac{C}{P\cdot\mathrm{MFU}(\text{shape},\text{batch})} \;+\; t_{\text{startup}} + t_{\text{eval}},
\]
and MFU (\cref{ch:pretrain}) is far from constant: narrow models launch small matmuls that cannot
fill the tensor cores, small batches leave the GPU waiting on kernel launches, and every run pays
a fixed startup cost (container start, data-loader warm-up, compilation). Consequences:
\begin{enumerate}
\item The tiny runs at the bottom of a study cost more device-time per FLOP than the big ones.
      Budget them in seconds, not FLOPs.
\item The \emph{wall-clock}-optimal model for a fixed number of device-hours can shift
      towards shapes with high MFU (wider, larger batch) relative to the FLOP-optimal model.
      The shift is largest where MFU still changes quickly with shape
      (\cref{fig:scalinglaws:wall}).
\item Measure MFU for each shape you might pick in the final run. Do not assume it.
\end{enumerate}
\end{systemslens}

% ======================================================================
\section{Kaplan versus Chinchilla}\label{sec:scalinglaws:kaplan}

Given $C$, how should we split it between $N$ and $D$? Two influential answers disagree.
Kaplan et al.\ (2020) found the loss-optimal model size grows as roughly $N_{\text{opt}}\propto
C^{0.73}$: most extra compute should go into a bigger model, trained on comparatively few
tokens. Hoffmann et al.\ (2022, ``Chinchilla'') found $N_{\text{opt}}\propto C^{0.5}$ and
$D_{\text{opt}}\propto C^{0.5}$: scale both equally, at about 20 tokens per parameter. The 70B
Chinchilla model, trained on 1.4T tokens for the same compute as the 280B Gopher model, beat it.

\begin{figure}[t]
\centering
\begin{tikzpicture}
\begin{loglogaxis}[width=0.62\linewidth, height=5cm, xmin=1e18, xmax=1e25,
  xlabel={compute $C$ (FLOPs)}, ylabel={$N_{\text{opt}}$ (relative)}, label style={font=\small},
  tick label style={font=\scriptsize},
  legend style={font=\scriptsize, draw=none, at={(1.02,0.5)}, anchor=west}]
\addplot[very thick, warnred, domain=1e18:1e25, samples=10] {(x/1e18)^0.73};
\addplot[very thick, sysblue, domain=1e18:1e25, samples=10] {(x/1e18)^0.5};
\legend{Kaplan: $C^{0.73}$, Chinchilla: $C^{0.50}$}
\end{loglogaxis}
\end{tikzpicture}
\caption{The two allocation rules, normalised to agree at $10^{18}$ FLOPs. Over seven decades of
compute they disagree by a factor of $10^{7\times0.23}\approx 40$ in model size. A small
difference in an exponent becomes a large difference in what you build.}
\label{fig:scalinglaws:kaplan}
\end{figure}

Later analyses attribute most of the gap to methodology rather than to the models (Porian et
al., 2024, and Pearce and Song, 2024, give careful reconciliations). The main suspects:
\begin{itemize}
  \item \textbf{Learning-rate schedules.} Kaplan's runs used one schedule length and read off
        intermediate losses for shorter token counts. A cosine schedule that has not finished
        decaying reports a loss well above what a run of that length would reach, which
        penalises long-data configurations. Chinchilla matched the schedule to each run's length.
        Our synthetic API reproduces this effect (\cref{sec:scalinglaws:wsd}).
  \item \textbf{Which $N$.} Kaplan counted non-embedding parameters and omitted the LM head from
        compute; at small scale that bends the curve.
  \item \textbf{Scale.} Kaplan's runs were smaller, where the corrections above matter most.
\end{itemize}
The lesson for practitioners is not ``Chinchilla is right''. It is that the allocation exponent
is a measured quantity whose value depends on choices you make in the study, so you must make
those choices deliberately and consistently.

% ======================================================================
\section{Three ways to fit an allocation law}\label{sec:scalinglaws:methods}

Hoffmann et al.\ estimated $N_{\text{opt}}(C)$ three ways. All three are worth knowing, because they
fail differently.

\subsection{Method 1: the lower envelope of training curves}\label{sec:scalinglaws:method1}

Train models of several sizes, each for a long run, and record loss against FLOPs consumed
along the way. For each compute value, the best model is whichever curve is lowest there: the
\emph{lower envelope} (\cref{fig:scalinglaws:envelope}). The envelope's contact points give
$N_{\text{opt}}(C)$. Method 1 is cheap (one run per size), but it needs the intermediate losses
to be valid stand-ins for completed runs, and with a cosine schedule they are not. A
warmup--stable--decay schedule fixes this (\cref{sec:scalinglaws:wsd}).

\begin{figure}[t]
\centering
\begin{tikzpicture}
\begin{axis}[xmode=log, width=0.78\linewidth, height=5.6cm, xmin=1e15, xmax=1e22, ymin=1.95, ymax=6,
  xlabel={compute $C$ (FLOPs)}, ylabel={loss}, label style={font=\small},
  tick label style={font=\scriptsize},
  legend style={font=\scriptsize, draw=none, at={(1.02,0.5)}, anchor=west}]
\addplot[thick, black!35] coordinates {(1e+15,6.958) (1.314e+15,6.75) (1.727e+15,6.559) (2.27e+15,6.382) (2.982e+15,6.218) (3.919e+15,6.067) (5.151e+15,5.927) (6.769e+15,5.798) (8.895e+15,5.679) (1.169e+16,5.569) (1.536e+16,5.467) (2.019e+16,5.373) (2.653e+16,5.286) (3.486e+16,5.206) (4.582e+16,5.132) (6.021e+16,5.064) (7.912e+16,5) (1.04e+17,4.942) (1.366e+17,4.888) (1.796e+17,4.838) (2.36e+17,4.792) (3.101e+17,4.75) (4.075e+17,4.71) (5.356e+17,4.674) (7.038e+17,4.64) (9.249e+17,4.609) (1.215e+18,4.581) (1.597e+18,4.554) (2.099e+18,4.53) (2.759e+18,4.507) (3.625e+18,4.486) (4.764e+18,4.467) (6.261e+18,4.449) (8.227e+18,4.433) (1.081e+19,4.418) (1.421e+19,4.404) (1.867e+19,4.391) (2.454e+19,4.379) (3.225e+19,4.368) (4.238e+19,4.357) (5.569e+19,4.348) (7.318e+19,4.339) (9.617e+19,4.331) (1.264e+20,4.324) (1.661e+20,4.317) (2.183e+20,4.311) (2.868e+20,4.305) (3.769e+20,4.299) (4.954e+20,4.294) (6.51e+20,4.29) (8.555e+20,4.285) (1.124e+21,4.281) (1.477e+21,4.278) (1.941e+21,4.274) (2.551e+21,4.271) (3.353e+21,4.268) (4.406e+21,4.266) (5.79e+21,4.263) (7.609e+21,4.261) (1e+22,4.259)};
\addplot[thick, black!35] coordinates {(1.314e+15,6.941) (1.727e+15,6.673) (2.27e+15,6.426) (2.982e+15,6.197) (3.919e+15,5.986) (5.151e+15,5.791) (6.769e+15,5.611) (8.895e+15,5.445) (1.169e+16,5.291) (1.536e+16,5.149) (2.019e+16,5.018) (2.653e+16,4.897) (3.486e+16,4.785) (4.582e+16,4.681) (6.021e+16,4.586) (7.912e+16,4.497) (1.04e+17,4.416) (1.366e+17,4.34) (1.796e+17,4.271) (2.36e+17,4.206) (3.101e+17,4.147) (4.075e+17,4.092) (5.356e+17,4.041) (7.038e+17,3.995) (9.249e+17,3.951) (1.215e+18,3.911) (1.597e+18,3.874) (2.099e+18,3.84) (2.759e+18,3.809) (3.625e+18,3.78) (4.764e+18,3.753) (6.261e+18,3.728) (8.227e+18,3.705) (1.081e+19,3.684) (1.421e+19,3.664) (1.867e+19,3.646) (2.454e+19,3.629) (3.225e+19,3.614) (4.238e+19,3.6) (5.569e+19,3.586) (7.318e+19,3.574) (9.617e+19,3.563) (1.264e+20,3.552) (1.661e+20,3.543) (2.183e+20,3.534) (2.868e+20,3.526) (3.769e+20,3.518) (4.954e+20,3.511) (6.51e+20,3.505) (8.555e+20,3.499) (1.124e+21,3.493) (1.477e+21,3.488) (1.941e+21,3.483) (2.551e+21,3.479) (3.353e+21,3.475) (4.406e+21,3.471) (5.79e+21,3.468) (7.609e+21,3.465) (1e+22,3.462)};
\addplot[thick, black!35] coordinates {(3.919e+15,6.462) (5.151e+15,6.189) (6.769e+15,5.938) (8.895e+15,5.705) (1.169e+16,5.491) (1.536e+16,5.292) (2.019e+16,5.109) (2.653e+16,4.94) (3.486e+16,4.784) (4.582e+16,4.639) (6.021e+16,4.506) (7.912e+16,4.382) (1.04e+17,4.268) (1.366e+17,4.163) (1.796e+17,4.066) (2.36e+17,3.976) (3.101e+17,3.893) (4.075e+17,3.817) (5.356e+17,3.746) (7.038e+17,3.68) (9.249e+17,3.62) (1.215e+18,3.564) (1.597e+18,3.513) (2.099e+18,3.465) (2.759e+18,3.421) (3.625e+18,3.38) (4.764e+18,3.343) (6.261e+18,3.308) (8.227e+18,3.276) (1.081e+19,3.246) (1.421e+19,3.219) (1.867e+19,3.194) (2.454e+19,3.17) (3.225e+19,3.149) (4.238e+19,3.129) (5.569e+19,3.11) (7.318e+19,3.093) (9.617e+19,3.078) (1.264e+20,3.063) (1.661e+20,3.05) (2.183e+20,3.037) (2.868e+20,3.026) (3.769e+20,3.015) (4.954e+20,3.005) (6.51e+20,2.996) (8.555e+20,2.988) (1.124e+21,2.98) (1.477e+21,2.973) (1.941e+21,2.967) (2.551e+21,2.961) (3.353e+21,2.955) (4.406e+21,2.95) (5.79e+21,2.945) (7.609e+21,2.941) (1e+22,2.937)};
\addplot[thick, black!35] coordinates {(1.169e+16,6.162) (1.536e+16,5.885) (2.019e+16,5.63) (2.653e+16,5.393) (3.486e+16,5.175) (4.582e+16,4.973) (6.021e+16,4.787) (7.912e+16,4.615) (1.04e+17,4.456) (1.366e+17,4.309) (1.796e+17,4.173) (2.36e+17,4.047) (3.101e+17,3.932) (4.075e+17,3.825) (5.356e+17,3.726) (7.038e+17,3.634) (9.249e+17,3.55) (1.215e+18,3.472) (1.597e+18,3.4) (2.099e+18,3.333) (2.759e+18,3.272) (3.625e+18,3.215) (4.764e+18,3.163) (6.261e+18,3.114) (8.227e+18,3.07) (1.081e+19,3.028) (1.421e+19,2.99) (1.867e+19,2.955) (2.454e+19,2.922) (3.225e+19,2.892) (4.238e+19,2.864) (5.569e+19,2.838) (7.318e+19,2.815) (9.617e+19,2.793) (1.264e+20,2.772) (1.661e+20,2.754) (2.183e+20,2.736) (2.868e+20,2.72) (3.769e+20,2.706) (4.954e+20,2.692) (6.51e+20,2.679) (8.555e+20,2.668) (1.124e+21,2.657) (1.477e+21,2.647) (1.941e+21,2.638) (2.551e+21,2.629) (3.353e+21,2.622) (4.406e+21,2.614) (5.79e+21,2.608) (7.609e+21,2.601) (1e+22,2.596)};
\addplot[thick, black!35] coordinates {(4.582e+16,5.702) (6.021e+16,5.442) (7.912e+16,5.201) (1.04e+17,4.979) (1.366e+17,4.774) (1.796e+17,4.585) (2.36e+17,4.409) (3.101e+17,4.248) (4.075e+17,4.098) (5.356e+17,3.96) (7.038e+17,3.833) (9.249e+17,3.715) (1.215e+18,3.606) (1.597e+18,3.505) (2.099e+18,3.412) (2.759e+18,3.327) (3.625e+18,3.247) (4.764e+18,3.174) (6.261e+18,3.106) (8.227e+18,3.044) (1.081e+19,2.986) (1.421e+19,2.933) (1.867e+19,2.883) (2.454e+19,2.838) (3.225e+19,2.796) (4.238e+19,2.757) (5.569e+19,2.721) (7.318e+19,2.688) (9.617e+19,2.657) (1.264e+20,2.629) (1.661e+20,2.603) (2.183e+20,2.579) (2.868e+20,2.556) (3.769e+20,2.536) (4.954e+20,2.517) (6.51e+20,2.499) (8.555e+20,2.483) (1.124e+21,2.468) (1.477e+21,2.454) (1.941e+21,2.441) (2.551e+21,2.429) (3.353e+21,2.418) (4.406e+21,2.408) (5.79e+21,2.399) (7.609e+21,2.39) (1e+22,2.382)};
\addplot[thick, black!35] coordinates {(1.366e+17,5.599) (1.796e+17,5.335) (2.36e+17,5.09) (3.101e+17,4.864) (4.075e+17,4.656) (5.356e+17,4.463) (7.038e+17,4.285) (9.249e+17,4.12) (1.215e+18,3.968) (1.597e+18,3.828) (2.099e+18,3.698) (2.759e+18,3.578) (3.625e+18,3.467) (4.764e+18,3.365) (6.261e+18,3.271) (8.227e+18,3.183) (1.081e+19,3.103) (1.421e+19,3.028) (1.867e+19,2.959) (2.454e+19,2.896) (3.225e+19,2.837) (4.238e+19,2.783) (5.569e+19,2.733) (7.318e+19,2.686) (9.617e+19,2.644) (1.264e+20,2.604) (1.661e+20,2.567) (2.183e+20,2.534) (2.868e+20,2.503) (3.769e+20,2.474) (4.954e+20,2.447) (6.51e+20,2.423) (8.555e+20,2.4) (1.124e+21,2.379) (1.477e+21,2.36) (1.941e+21,2.342) (2.551e+21,2.325) (3.353e+21,2.31) (4.406e+21,2.296) (5.79e+21,2.283) (7.609e+21,2.271) (1e+22,2.259)};
\addplot[thick, black!35] coordinates {(4.075e+17,5.551) (5.356e+17,5.282) (7.038e+17,5.033) (9.249e+17,4.804) (1.215e+18,4.591) (1.597e+18,4.395) (2.099e+18,4.214) (2.759e+18,4.047) (3.625e+18,3.892) (4.764e+18,3.749) (6.261e+18,3.617) (8.227e+18,3.495) (1.081e+19,3.383) (1.421e+19,3.279) (1.867e+19,3.183) (2.454e+19,3.094) (3.225e+19,3.012) (4.238e+19,2.936) (5.569e+19,2.866) (7.318e+19,2.801) (9.617e+19,2.742) (1.264e+20,2.686) (1.661e+20,2.635) (2.183e+20,2.588) (2.868e+20,2.545) (3.769e+20,2.505) (4.954e+20,2.467) (6.51e+20,2.433) (8.555e+20,2.401) (1.124e+21,2.372) (1.477e+21,2.345) (1.941e+21,2.32) (2.551e+21,2.297) (3.353e+21,2.276) (4.406e+21,2.256) (5.79e+21,2.238) (7.609e+21,2.221) (1e+22,2.205)};
\addplot[very thick, warnred, dashed] coordinates {(1e+15,6.954) (1.512e+15,6.627) (2.285e+15,6.321) (3.455e+15,6.034) (5.223e+15,5.766) (7.897e+15,5.514) (1.194e+16,5.278) (1.805e+16,5.058) (2.728e+16,4.851) (4.125e+16,4.657) (6.236e+16,4.476) (9.427e+16,4.306) (1.425e+17,4.147) (2.154e+17,3.998) (3.257e+17,3.859) (4.924e+17,3.728) (7.444e+17,3.606) (1.125e+18,3.491) (1.701e+18,3.384) (2.572e+18,3.284) (3.888e+18,3.189) (5.878e+18,3.101) (8.886e+18,3.019) (1.343e+19,2.941) (2.031e+19,2.869) (3.07e+19,2.801) (4.642e+19,2.738) (7.017e+19,2.678) (1.061e+20,2.622) (1.604e+20,2.57) (2.424e+20,2.521) (3.665e+20,2.476) (5.541e+20,2.433) (8.377e+20,2.393) (1.266e+21,2.355) (1.914e+21,2.32) (2.894e+21,2.287) (4.375e+21,2.256) (6.615e+21,2.227) (1e+22,2.2)};
\end{axis}
\end{tikzpicture}
\caption{Method 1. Each grey curve is one model size ($3\times10^6$ to $3\times10^9$ non-embedding
parameters) trained for longer and longer; the dashed red line is their lower envelope. Data:
the noise-free part of the synthetic ground truth in \code{llmsys.scaling}, not a real model.}
\label{fig:scalinglaws:envelope}
\end{figure}

\subsection{Method 2: IsoFLOP profiles}\label{sec:scalinglaws:method2}

Pick a handful of compute budgets $C_1<C_2<\dots$. For each budget, train several model sizes
$N_{ij}$, each on exactly $D_{ij}=C_i/(6N_{ij})$ tokens with a schedule matched to that length,
and record the final loss $L_{ij}$. Plotted against $\log N$, each budget's losses form a valley:
too small a model cannot use the data, too large a model cannot take enough steps
(\cref{fig:scalinglaws:isoflop}). The bottom of each valley is $N_{\text{opt}}(C_i)$; a power law
through those points is the allocation law.

\begin{figure}[t]
\centering
\begin{tikzpicture}
\begin{groupplot}[group style={group size=2 by 1, horizontal sep=1.6cm},
  width=0.5\linewidth, height=5.2cm, label style={font=\small},
  tick label style={font=\scriptsize}]
\nextgroupplot[xmode=log, xlabel={non-embedding params $N$}, ylabel={final loss},
  xmin=3e6, xmax=3e10, ymin=2.2, ymax=4.0]
\addplot[thick, sysblue] coordinates {(2.406e+06,5.101) (2.617e+06,5.041) (2.847e+06,4.984) (3.097e+06,4.93) (3.369e+06,4.878) (3.666e+06,4.828) (3.988e+06,4.781) (4.338e+06,4.736) (4.719e+06,4.694) (5.134e+06,4.654) (5.585e+06,4.616) (6.076e+06,4.58) (6.61e+06,4.547) (7.191e+06,4.515) (7.823e+06,4.486) (8.511e+06,4.459) (9.259e+06,4.434) (1.007e+07,4.411) (1.096e+07,4.39) (1.192e+07,4.371) (1.297e+07,4.353) (1.411e+07,4.338) (1.535e+07,4.325) (1.67e+07,4.313) (1.816e+07,4.304) (1.976e+07,4.296) (2.15e+07,4.29) (2.339e+07,4.286) (2.544e+07,4.284) (2.768e+07,4.283) (3.011e+07,4.284) (3.276e+07,4.287) (3.563e+07,4.292) (3.877e+07,4.299) (4.217e+07,4.307) (4.588e+07,4.317) (4.991e+07,4.329) (5.43e+07,4.343) (5.907e+07,4.358) (6.426e+07,4.375) (6.991e+07,4.394) (7.605e+07,4.415) (8.274e+07,4.437) (9.001e+07,4.461) (9.792e+07,4.488) (1.065e+08,4.515) (1.159e+08,4.545) (1.261e+08,4.576) (1.371e+08,4.61) (1.492e+08,4.645) (1.623e+08,4.682) (1.766e+08,4.721) (1.921e+08,4.762) (2.09e+08,4.804) (2.273e+08,4.849) (2.473e+08,4.896) (2.691e+08,4.944) (2.927e+08,4.995) (3.184e+08,5.047) (3.464e+08,5.102)};
\addplot[only marks, mark=o, mark size=1.4pt, sysblue] coordinates {(3.608e+06,4.849) (7.217e+06,4.499) (1.443e+07,4.337) (2.887e+07,4.28) (5.774e+07,4.351) (1.155e+08,4.542) (2.309e+08,4.845)};
\addplot[only marks, mark=star, mark size=3pt, sysblue] coordinates {(2.768e+07,4.283)};
\node[font=\scriptsize, text=sysblue] at (axis cs:2.768e+07,4.223) {$10^{17}$};
\addplot[thick, keygreen] coordinates {(7.607e+06,4.038) (8.276e+06,3.999) (9.003e+06,3.962) (9.794e+06,3.926) (1.066e+07,3.892) (1.159e+07,3.86) (1.261e+07,3.83) (1.372e+07,3.801) (1.492e+07,3.773) (1.624e+07,3.747) (1.766e+07,3.723) (1.921e+07,3.7) (2.09e+07,3.679) (2.274e+07,3.659) (2.474e+07,3.64) (2.691e+07,3.623) (2.928e+07,3.608) (3.185e+07,3.594) (3.465e+07,3.581) (3.77e+07,3.569) (4.101e+07,3.559) (4.461e+07,3.55) (4.853e+07,3.542) (5.28e+07,3.536) (5.744e+07,3.531) (6.249e+07,3.527) (6.798e+07,3.525) (7.395e+07,3.524) (8.045e+07,3.524) (8.752e+07,3.525) (9.521e+07,3.527) (1.036e+08,3.531) (1.127e+08,3.536) (1.226e+08,3.542) (1.334e+08,3.549) (1.451e+08,3.558) (1.578e+08,3.568) (1.717e+08,3.579) (1.868e+08,3.591) (2.032e+08,3.604) (2.211e+08,3.619) (2.405e+08,3.635) (2.616e+08,3.652) (2.846e+08,3.671) (3.096e+08,3.69) (3.369e+08,3.711) (3.665e+08,3.733) (3.987e+08,3.757) (4.337e+08,3.782) (4.718e+08,3.808) (5.133e+08,3.835) (5.584e+08,3.864) (6.075e+08,3.894) (6.608e+08,3.925) (7.189e+08,3.958) (7.821e+08,3.992) (8.508e+08,4.028) (9.256e+08,4.064) (1.007e+09,4.103) (1.095e+09,4.143)};
\addplot[only marks, mark=o, mark size=1.4pt, keygreen] coordinates {(1.141e+07,3.865) (2.282e+07,3.653) (4.564e+07,3.568) (9.129e+07,3.527) (1.826e+08,3.585) (3.651e+08,3.731) (7.303e+08,3.96)};
\addplot[only marks, mark=star, mark size=3pt, keygreen] coordinates {(8.045e+07,3.524)};
\node[font=\scriptsize, text=keygreen] at (axis cs:8.045e+07,3.464) {$10^{18}$};
\addplot[thick, napkin] coordinates {(2.406e+07,3.319) (2.617e+07,3.293) (2.847e+07,3.269) (3.097e+07,3.246) (3.369e+07,3.224) (3.666e+07,3.203) (3.988e+07,3.183) (4.338e+07,3.164) (4.719e+07,3.147) (5.134e+07,3.13) (5.585e+07,3.115) (6.076e+07,3.1) (6.61e+07,3.087) (7.191e+07,3.074) (7.823e+07,3.062) (8.511e+07,3.052) (9.259e+07,3.042) (1.007e+08,3.034) (1.096e+08,3.026) (1.192e+08,3.019) (1.297e+08,3.013) (1.411e+08,3.008) (1.535e+08,3.004) (1.67e+08,3.001) (1.816e+08,2.998) (1.976e+08,2.997) (2.15e+08,2.996) (2.339e+08,2.996) (2.544e+08,2.998) (2.768e+08,2.999) (3.011e+08,3.002) (3.276e+08,3.006) (3.563e+08,3.01) (3.877e+08,3.016) (4.217e+08,3.022) (4.588e+08,3.029) (4.991e+08,3.037) (5.43e+08,3.046) (5.907e+08,3.055) (6.426e+08,3.066) (6.991e+08,3.077) (7.605e+08,3.089) (8.274e+08,3.102) (9.001e+08,3.116) (9.792e+08,3.131) (1.065e+09,3.146) (1.159e+09,3.163) (1.261e+09,3.18) (1.371e+09,3.199) (1.492e+09,3.218) (1.623e+09,3.238) (1.766e+09,3.259) (1.921e+09,3.281) (2.09e+09,3.304) (2.273e+09,3.328) (2.473e+09,3.353) (2.691e+09,3.379) (2.927e+09,3.406) (3.184e+09,3.434) (3.464e+09,3.463)};
\addplot[only marks, mark=o, mark size=1.4pt, napkin] coordinates {(3.608e+07,3.2) (7.217e+07,3.071) (1.443e+08,3.01) (2.887e+08,2.999) (5.774e+08,3.058) (1.155e+09,3.161) (2.309e+09,3.333)};
\addplot[only marks, mark=star, mark size=3pt, napkin] coordinates {(2.15e+08,2.996)};
\node[font=\scriptsize, text=napkin] at (axis cs:2.15e+08,2.936) {$10^{19}$};
\addplot[thick, labpurple] coordinates {(7.607e+07,2.832) (8.276e+07,2.815) (9.003e+07,2.799) (9.794e+07,2.784) (1.066e+08,2.77) (1.159e+08,2.756) (1.261e+08,2.743) (1.372e+08,2.731) (1.492e+08,2.72) (1.624e+08,2.71) (1.766e+08,2.7) (1.921e+08,2.69) (2.09e+08,2.682) (2.274e+08,2.674) (2.474e+08,2.667) (2.691e+08,2.66) (2.928e+08,2.655) (3.185e+08,2.649) (3.465e+08,2.645) (3.77e+08,2.641) (4.101e+08,2.638) (4.461e+08,2.635) (4.853e+08,2.633) (5.28e+08,2.631) (5.744e+08,2.63) (6.249e+08,2.63) (6.798e+08,2.631) (7.395e+08,2.631) (8.045e+08,2.633) (8.752e+08,2.635) (9.521e+08,2.638) (1.036e+09,2.641) (1.127e+09,2.645) (1.226e+09,2.649) (1.334e+09,2.654) (1.451e+09,2.66) (1.578e+09,2.666) (1.717e+09,2.673) (1.868e+09,2.681) (2.032e+09,2.689) (2.211e+09,2.697) (2.405e+09,2.706) (2.616e+09,2.716) (2.846e+09,2.727) (3.096e+09,2.738) (3.369e+09,2.749) (3.665e+09,2.762) (3.987e+09,2.775) (4.337e+09,2.788) (4.718e+09,2.802) (5.133e+09,2.817) (5.584e+09,2.833) (6.075e+09,2.849) (6.608e+09,2.866) (7.189e+09,2.883) (7.821e+09,2.901) (8.508e+09,2.92) (9.256e+09,2.94) (1.007e+10,2.96) (1.095e+10,2.981)};
\addplot[only marks, mark=o, mark size=1.4pt, labpurple] coordinates {(1.141e+08,2.768) (2.282e+08,2.677) (4.564e+08,2.631) (9.129e+08,2.635) (1.826e+09,2.682) (3.651e+09,2.773) (7.303e+09,2.885)};
\addplot[only marks, mark=star, mark size=3pt, labpurple] coordinates {(6.249e+08,2.63)};
\node[font=\scriptsize, text=labpurple] at (axis cs:6.249e+08,2.570) {$10^{20}$};

\nextgroupplot[xmode=log, ymode=log, xlabel={compute $C$ (FLOPs)}, ylabel={$N_{\text{opt}}$},
  xmin=1e16, xmax=1e23]
\addplot[only marks, mark=*, sysblue] coordinates {(1e+17,2.768e+07) (1e+18,8.045e+07) (1e+19,2.15e+08) (1e+20,6.249e+08)};
\addplot[thick, warnred, dashed] coordinates {(1e+16,9.931e+06) (3.162e+16,1.665e+07) (1e+17,2.791e+07) (3.162e+17,4.679e+07) (1e+18,7.844e+07) (3.162e+18,1.315e+08) (1e+19,2.205e+08) (3.162e+19,3.696e+08) (1e+20,6.196e+08) (3.162e+20,1.039e+09) (1e+21,1.741e+09) (3.162e+21,2.919e+09) (1e+22,4.894e+09) (3.162e+22,8.205e+09) (1e+23,1.376e+10)};
\node[font=\scriptsize, anchor=north west, align=left] at (rel axis cs:0.03,0.97)
  {fit: $N_{\text{opt}}\propto C^{0.449}$\\truth: $C^{0.453}$};
\end{groupplot}
\end{tikzpicture}
\caption{Method 2 on the synthetic ground truth. Left: four IsoFLOP profiles
($10^{17}$--$10^{20}$ FLOPs); dots are noisy ``runs'', lines the noise-free valley, stars the
minima. Right: the minima and the fitted power law, extrapolated three decades. The fitted
exponent recovers the hidden $\beta/(\alpha+\beta)=0.453$ closely because the grid brackets every
minimum.}
\label{fig:scalinglaws:isoflop}
\end{figure}

There are two ways to read off the minimum. Taking the best run is robust but quantised to the
sizes you tried. With sizes spaced by a factor $r$, the true optimum can be up to $\sqrt r$ away.
Fitting a parabola in $\log N$ and taking its vertex interpolates between grid points, but it
trusts the parabola, which is only a local approximation to the valley. The vertex is only
meaningful if the grid brackets it: if the best run sits at the edge of the grid, the budget's
grid was mis-centred, and the right response is to add runs, not to extrapolate the parabola.
\Cref{lst:scalinglaws:isoflop} implements both readings with that guard.

\begin{lstlisting}[caption={IsoFLOP minima with a bracketing check (abridged from \code{llmsys/scaling.py}).},label={lst:scalinglaws:isoflop}]
def isoflop_optima(runs, method="min", rel_tol=0.05):
    out = []
    for c, group in group_by_budget(runs, rel_tol).items():  # budgets drift
        best = min(group, key=lambda r: r.loss)
        n_opt, loss, used = best.params, best.loss, "min"
        if method == "quadratic" and len(group) >= 3:
            ln = np.log([g.params for g in group]); mu = ln.mean()
            c2, c1, c0 = np.polyfit(ln - mu, [g.loss for g in group], 2)
            v = mu - c1 / (2 * c2)
            if c2 > 0 and ln.min() <= v <= ln.max():   # convex, bracketed
                n_opt, loss = math.exp(v), c0 - c1**2 / (4*c2)
                used = "quadratic"
        ...
\end{lstlisting}

\begin{pitfall}[Budgets are never exact]
Token counts must be whole numbers of batches, and model sizes come in multiples of the head
dimension, so the ``$10^{18}$-FLOP'' runs actually spend $0.98$--$1.02\times10^{18}$. Group runs by
\emph{relative} tolerance, and use each group's mean compute when you fit the power law. Never use
the nominal budget. A 2\% bias in $C$ is invisible on a plot but shifts the intercept of every
extrapolation.
\end{pitfall}

\subsection{Method 3: a parametric loss surface}\label{sec:scalinglaws:method3}

Instead of locating minima, fit a formula for the whole surface,
\begin{equation}
\hat L(N,D) \;=\; E + \frac{A}{N^{\alpha}} + \frac{B}{D^{\beta}},
\label{eq:scalinglaws:chin}
\end{equation}
to \emph{every} run, and then minimise it analytically. Minimising $\hat L$ subject to $6ND = C$
with a Lagrange multiplier gives
\begin{equation}
N_{\text{opt}}(C) = G\Bigl(\frac{C}{6}\Bigr)^{\frac{\beta}{\alpha+\beta}},\qquad
D_{\text{opt}}(C) = G^{-1}\Bigl(\frac{C}{6}\Bigr)^{\frac{\alpha}{\alpha+\beta}},\qquad
G = \Bigl(\frac{\alpha A}{\beta B}\Bigr)^{\frac{1}{\alpha+\beta}}.
\label{eq:scalinglaws:closed}
\end{equation}
The derivation is three lines (\cref{ex:scalinglaws:lagrange}). The exponents satisfy
$a+b=1$, as they must when $C\propto ND$. Equal exponents $\alpha=\beta$ give the Chinchilla
$C^{0.5}$ rule; a smaller data exponent $\beta$ pushes compute towards data.

The fit itself is a small but delicate optimisation problem. Hoffmann et al.\ minimised a Huber
loss on the \emph{log} of the loss, with the surface written as a log-sum-exp so that every
parameter is unconstrained:
\begin{equation}
\min_{a,b,e,\alpha,\beta}\;\sum_i \mathrm{Huber}_\delta\Bigl(
\mathrm{LSE}\bigl(a-\alpha\log N_i,\; b-\beta\log D_i,\; e\bigr) - \log L_i\Bigr),
\qquad A=e^a,\;B=e^b,\;E=e^e,
\end{equation}
with $\delta=10^{-3}$, from a grid of initial points, keeping the best. The Huber loss down-weights
outlier runs, such as one with a loss spike. Working in log space treats relative errors equally at
all scales. The grid of starts is not optional: with few runs the objective is nearly flat along
directions that trade $E$ against $A$ and $B$, and local optimisers stop in very different places.

\begin{lstlisting}[caption={Huber/LSE fit of the parametric law from a grid of starts (from \code{llmsys/scaling.py}).},label={lst:scalinglaws:chin}]
def fit_chinchilla(n, d, loss, delta=1e-3, inits=None, max_iter=100):
    ln, ld, ll = (torch.tensor(np.log(v), dtype=torch.float64)[None]
                  for v in (n, d, loss))
    inits = inits or [(a0, b0, e0, al, be) for a0 in (0., 5., 10.)
                      for b0 in (0., 5., 10.) for e0 in (-1., .5)
                      for al in (.1, .4) for be in (.1, .4)]          # 72 starts
    best, best_obj = None, float("inf")
    for init in inits:
        q = torch.tensor([init], dtype=torch.float64, requires_grad=True)
        opt = torch.optim.LBFGS([q], max_iter=max_iter,
                                line_search_fn="strong_wolfe")
        def closure():
            opt.zero_grad()
            a, b, e, al, be = (q[:, i:i+1] for i in range(5))
            pred = torch.logsumexp(torch.stack(
                [a - al*ln, b - be*ld, e.expand(-1, ln.shape[1])]), 0)
            obj = huber(pred - ll, delta).sum()
            obj.backward(); return obj
        opt.step(closure)
        ...   # keep the start with the lowest objective and alpha, beta > 0
\end{lstlisting}

\begin{deepdive}[When the three methods disagree]
In the Chinchilla paper, methods 1 and 2 agreed ($a\approx0.50$) and method 3 gave a somewhat
different exponent. A later replication attempt (Besiroglu et al., 2024) reconstructed the data
from the paper's figures and reported that the published method-3 parameters fit the data
poorly and imply implausibly narrow confidence intervals. Their refit moved method 3 into
agreement with the other two. The practical moral: report all the methods you can afford, look
at the residuals, and bootstrap (\cref{sec:scalinglaws:uncertainty}) before believing a third
significant digit of anything.
\end{deepdive}

% ======================================================================
\section{Hyperparameters across scale}\label{sec:scalinglaws:hparams}

A scaling law describes \emph{well-tuned} models. If every size uses the same learning rate,
the small models are under-stepped or the large ones over-stepped, and the measured curve
mixes the scaling of the architecture with the mis-tuning of the optimizer. Modern recipes
handle this in one of two ways: fit scaling laws for the hyperparameters themselves, or
reparameterise the model so that the optimal hyperparameters stop depending on size.

\subsection{Learning rate}\label{sec:scalinglaws:lr}

For a fixed size and token budget, final loss as a function of $\log(\text{LR})$ is a smooth,
roughly parabolic basin with a cliff on the high side, where training diverges
(\cref{fig:scalinglaws:basin}). Its minimum moves to smaller learning rates as models get wider.
The procedure is direct: sweep the LR at three or four small sizes, take each basin's vertex
(a quadratic in $\log\text{LR}$, ignoring diverged runs), and fit $\text{lr}^*(N)=kN^{a}$. DeepSeek's
LLM report (2024) is a public example of fitting both LR and batch size as power laws in compute
this way, and a large 2025 hyperparameter study from StepFun mapped the LR--batch landscape by
brute-force grids and found it convex with clean minimisers.

\begin{figure}[t]
\centering
\begin{tikzpicture}
\begin{axis}[xmode=log, width=0.72\linewidth, height=5.2cm, xmin=2e-4, xmax=2e-2,
  xlabel={peak learning rate}, ylabel={final loss (20 tokens/param)}, label style={font=\small},
  tick label style={font=\scriptsize},
  legend style={font=\scriptsize, draw=none, at={(1.02,0.5)}, anchor=west}]
\addplot[very thick, sysblue] coordinates {(0.0002,6.206) (0.0002251,6.158) (0.0002533,6.112) (0.000285,6.067) (0.0003207,6.024) (0.0003609,5.983) (0.0004062,5.944) (0.0004571,5.906) (0.0005144,5.87) (0.0005789,5.835) (0.0006514,5.802) (0.000733,5.771) (0.0008249,5.742) (0.0009283,5.714) (0.001045,5.688) (0.001176,5.664) (0.001323,5.641) (0.001489,5.62) (0.001675,5.6) (0.001885,5.583) (0.002122,5.567) (0.002388,5.552) (0.002687,5.54) (0.003024,5.528) (0.003403,5.519) (0.003829,5.511) (0.004309,5.505) (0.004849,5.501) (0.005457,5.498) (0.006141,5.497) (0.00691,5.498) (0.007776,5.5) (0.008751,5.504) (0.009848,5.51) (0.01108,5.518) (0.01247,5.527) (0.01403,5.537) (0.01579,5.55) (0.01777,5.564) (0.02,5.579)};
\addlegendentry{$N\approx10^{7}$}
\draw[sysblue, dashed] (axis cs:0.00621,2.0) -- (axis cs:0.00621,6);
\addplot[very thick, keygreen] coordinates {(0.0002,3.933) (0.0002251,3.896) (0.0002533,3.86) (0.000285,3.827) (0.0003207,3.794) (0.0003609,3.764) (0.0004062,3.735) (0.0004571,3.707) (0.0005144,3.682) (0.0005789,3.658) (0.0006514,3.636) (0.000733,3.615) (0.0008249,3.596) (0.0009283,3.579) (0.001045,3.563) (0.001176,3.549) (0.001323,3.537) (0.001489,3.527) (0.001675,3.518) (0.001885,3.511) (0.002122,3.505) (0.002388,3.501) (0.002687,3.499) (0.003024,3.499) (0.003403,3.5) (0.003829,3.503) (0.004309,3.507) (0.004849,3.513) (0.005457,3.521) (0.006141,3.531) (0.00691,3.542) (0.007776,3.555) (0.008751,3.569) (0.009848,3.586) (0.01108,3.603) (0.01247,3.623) (0.01403,3.644) (0.01579,3.667)};
\addlegendentry{$N\approx10^{8}$}
\draw[keygreen, dashed] (axis cs:0.002954,2.0) -- (axis cs:0.002954,6);
\addplot[very thick, napkin] coordinates {(0.0002,2.867) (0.0002251,2.839) (0.0002533,2.813) (0.000285,2.789) (0.0003207,2.766) (0.0003609,2.745) (0.0004062,2.726) (0.0004571,2.708) (0.0005144,2.692) (0.0005789,2.677) (0.0006514,2.665) (0.000733,2.654) (0.0008249,2.644) (0.0009283,2.636) (0.001045,2.63) (0.001176,2.626) (0.001323,2.623) (0.001489,2.622) (0.001675,2.623) (0.001885,2.625) (0.002122,2.629) (0.002388,2.635) (0.002687,2.642) (0.003024,2.651) (0.003403,2.662) (0.003829,2.674) (0.004309,2.688) (0.004849,2.704) (0.005457,2.721) (0.006141,2.741) (0.00691,2.761) (0.007776,2.784) (0.008751,2.808)};
\addlegendentry{$N\approx10^{9}$}
\draw[napkin, dashed] (axis cs:0.001508,2.0) -- (axis cs:0.001508,6);

\end{axis}
\end{tikzpicture}
\caption{Learning-rate basins on the synthetic ground truth for three model sizes, each trained
on 20 tokens per parameter. Vertical lines mark each optimum; it drifts down as the model
grows. Curves end where training diverges. A study that uses one LR for all sizes reads the
\emph{wrong} point of each basin.}
\label{fig:scalinglaws:basin}
\end{figure}

\subsection{Batch size and the critical batch}\label{sec:scalinglaws:batch}

Doubling the batch halves the number of optimizer steps but, past some point, does not halve
the number of steps \emph{needed}. McCandlish et al.\ (2018) model the trade-off between the
optimizer steps $S$ and examples $E$ needed to reach a target loss as
\begin{equation}
\Bigl(\frac{S}{S_{\min}}-1\Bigr)\Bigl(\frac{E}{E_{\min}}-1\Bigr) = 1,
\qquad B_{\text{crit}} = \frac{E_{\min}}{S_{\min}},
\label{eq:scalinglaws:critbatch}
\end{equation}
a hyperbola (\cref{fig:scalinglaws:critbatch}). Below $B_{\text{crit}}$, larger batches are
almost free: the same data, fewer steps. Above it, each doubling of the batch wastes most of
the extra data. At $B=B_{\text{crit}}$ you pay exactly twice the minimum of both. The critical
batch is approximately the \emph{gradient noise scale}: the ratio of the trace of the per-example
gradient covariance to the squared norm of the mean gradient. It grows as the loss falls, so
bigger and longer runs can use larger batches. That is why large-scale schedules ramp the
batch size up during training.

\begin{figure}[t]
\centering
\begin{tikzpicture}
\begin{loglogaxis}[width=0.6\linewidth, height=5cm, xmin=1, xmax=100, ymin=1, ymax=100,
  xlabel={$S/S_{\min}$ (optimizer steps)}, ylabel={$E/E_{\min}$ (examples)},
  label style={font=\small}, tick label style={font=\scriptsize}]
\addplot[very thick, sysblue, domain=1.01:100, samples=200] {1 + 1/(x-1)};
\addplot[only marks, mark=*, mark size=2.2pt, warnred] coordinates {(2,2)};
\node[font=\scriptsize, anchor=south west] at (axis cs:2.1,2.1) {$B=B_{\text{crit}}$};
\node[font=\scriptsize, align=center, anchor=south] at (axis cs:30,1.25) {small batch: few wasted\\examples, many serial steps};
\node[font=\scriptsize, align=left, anchor=west] at (axis cs:2.6,30) {large batch: few steps,\\much wasted data};
\end{loglogaxis}
\end{tikzpicture}
\caption{The steps--examples trade-off of \cref{eq:scalinglaws:critbatch} (analytic). The
critical batch is the knee, where you pay $2\times$ the minimum of each. Systems want to sit at
or slightly above it: fewer, larger steps parallelise across more devices (\cref{ch:distributed}).}
\label{fig:scalinglaws:critbatch}
\end{figure}

\begin{systemslens}[Batch size is where the optimizer meets the cluster]
Data parallelism needs a large global batch: with 1{,}024 GPUs and 4 sequences of 4{,}096 tokens
each, one step already consumes 16.8M tokens. The critical batch says whether that is
efficient. If $B_{\text{crit}}$ is 4M tokens at your loss, a 16M-token batch wastes most of
the extra data, and the honest choice is fewer GPUs per run or more model parallelism
(\cref{ch:distributed,ch:sharding}). A scaling study that sweeps batch size at small scale tells
you how far data parallelism can be pushed at large scale.
\end{systemslens}

\subsection{\texorpdfstring{$\mu$P}{muP}: make the optimum stop moving}\label{sec:scalinglaws:mup}

The maximal-update parameterisation ($\mu$P; Yang et al., 2022) takes the other route. It
chooses per-layer initialisation scales and learning-rate multipliers so that, as width $n$ grows,
(A1) every activation coordinate stays $\Theta(1)$ at initialisation, and (A2) one optimizer
step changes every activation coordinate by $\Theta(1)$, no more and no less. For a hidden
weight matrix $W\in\R^{n\times n}$ the activation condition gives the familiar
$\sigma\propto1/\sqrt{n}$ initialisation. The update condition is the new part. With Adam,
the update $\Delta W$ has entries of size about $\eta$, its product with an activation
vector of $\Theta(1)$ coordinates sums $n$ correlated terms, and so the output changes by
$\Theta(\eta n)$. Keeping that $\Theta(1)$ requires $\eta\propto 1/n$ for hidden layers, while
the input embedding keeps a width-independent rate and the output layer's logits are scaled down
by $1/n$. Under these rules the LR basin of \cref{fig:scalinglaws:basin} stays put as width
grows, so you tune once on a narrow proxy model and transfer (``$\mu$Transfer'').

\begin{keyidea}[Two public recipes]
Two well-documented public recipes make a useful contrast. \textbf{MiniCPM} (2024) uses
$\mu$P to make hyperparameters width-invariant, fixes the aspect ratio, and fits the
loss surface (method 3) with cheap WSD branches. \textbf{DeepSeek LLM} (2024) uses standard
parameterisation, fits optimal LR and batch size as power laws in compute, and sizes the model
with IsoFLOP profiles (method 2). Both work; what they share is that \emph{every}
scale-dependent knob is either modelled or parameterised away.
\end{keyidea}

% ======================================================================
\section{Making studies cheap: warmup--stable--decay}\label{sec:scalinglaws:wsd}

Method 2 needs a separate run for every (size, token count) pair, because a cosine schedule's
loss at 50\% of training is not the loss a run of half the length would reach: the LR has not
decayed yet. Fitting $m$ sizes at $k$ durations therefore costs $mk$ full runs, and if you want
the durations to cover the same range as the sizes, the cost of a study grows like $n^2$.

The \emph{warmup--stable--decay} (WSD) schedule (used by MiniCPM, and as a multi-step variant
by DeepSeek) removes the coupling. Warm up, hold the LR constant for most of training, and decay
quickly over the last $\sim$10--20\% of steps. Because the stable phase does not depend on the
final length, you can save a checkpoint at any point of it and \emph{branch}: run a short decay
from that checkpoint and get, for the price of the decay, a loss that matches a full run of that
length closely (\cref{fig:scalinglaws:wsd}). One long stable run per size plus a few cheap
decays replaces $k$ separate runs.

\begin{figure}[t]
\centering
\begin{tikzpicture}
\begin{axis}[width=0.48\linewidth, height=4.6cm, xmin=0, xmax=1, ymin=0, ymax=1.08,
  xlabel={training progress}, ylabel={LR / peak}, label style={font=\small},
  tick label style={font=\scriptsize}, title={\small schedules}, title style={yshift=-4pt},
  legend style={font=\scriptsize, draw=none, at={(0.5,-0.42)}, anchor=north, legend columns=2}]
\addplot[very thick, sysblue, domain=0:1, samples=100] {x<0.05 ? x/0.05 : 0.1+0.9*0.5*(1+cos(deg(pi*(x-0.05)/0.95)))};
\addplot[very thick, napkin, domain=0:1, samples=200] {x<0.05 ? x/0.05 : (x<0.85 ? 1 : 1-0.9*(x-0.85)/0.15)};
\legend{cosine, WSD}
\end{axis}
\end{tikzpicture}\hfill
\begin{tikzpicture}[x=5.2cm, y=1cm, font=\scriptsize]
\draw[->] (0,0) -- (1.05,0) node[right] {steps};
\draw[very thick, napkin] (0,0.2) -- (0.05,1.4) -- (0.95,1.4);
\foreach \b/\lab in {0.3/{$D_1$}, 0.55/{$D_2$}, 0.8/{$D_3$}} {
  \fill[napkin] (\b,1.4) circle (1.4pt);
  \draw[very thick, warnred] (\b,1.4) -- (\b+0.1,0.35);
  \node[below] at (\b+0.1,0.3) {\lab};
}
\node[anchor=west, align=left] at (0.02,2.1) {one stable run (orange); checkpoints};
\node[anchor=west, align=left] at (0.02,1.75) {branch into short decays (red)};
\node[anchor=north, align=center, text width=5cm] at (0.5,-0.25)
  {three valid end-of-run losses for the cost of one run plus three decays};
\end{tikzpicture}
\caption{Left: cosine versus WSD learning-rate schedules (analytic). Right: branching decays off
a WSD stable phase gives a loss at several data sizes from one run.}
\label{fig:scalinglaws:wsd}
\end{figure}

\begin{napkinmath}[What branching saves]
Suppose you want losses for one model at token counts $D, 2D, 4D, 8D$. With cosine you need
four runs: $15D$ tokens. With WSD and decays of 10\% of each length, you need one $8D$ run plus
decays of $0.1D+0.2D+0.4D+0.8D=1.5D$, so $9.5D$ tokens in total, a 37\% saving. With many
log-spaced durations the saving approaches 40\%, and the stable run doubles as the
method-1 curve. The cost is infrastructure: you must checkpoint the full optimizer state at
every branch point, restore it exactly, and track which branch produced which loss, which is
exactly the bookkeeping of \cref{sec:pretrain:save} and \cref{sec:scalinglaws:running}.
\end{napkinmath}

The synthetic API of this chapter models the cosine effect. Each run reports losses at
$k=1..n_{\text{evals}}$ evenly spaced points, and the early ones sit above the loss of a finished
run of that length by an amount proportional to the not-yet-decayed LR. Reading those
intermediate numbers as completed runs biases the fitted data exponent, which is the Kaplan
pitfall in miniature (\cref{ex:scalinglaws:midrun}).

% ======================================================================
\section{Beyond the basic law}\label{sec:scalinglaws:beyond}

\paragraph{Finite data and repetition.} Chinchilla assumes fresh tokens. When unique data runs
out, repeated epochs are worth less. Muennighoff et al.\ (2023) fit an ``effective data''
model in which the value of repeated tokens decays with the number of repetitions. Up to about
four epochs the loss is close to what the same number of unique tokens would give, and beyond that
the returns fall off sharply. Data-constrained regimes push the optimum towards larger models
trained on fewer epochs.

\paragraph{Data mixtures.} The proportions of web, code, books and so on can themselves be
treated as scaling-law inputs. Fitting mixture laws on small proxies is an active research area
and a hard one: small-scale rankings of mixtures do not always survive scaling, and simply
picking the best mixture at small scale is a strong baseline.
\Cref{ch:data} covers how the data itself is built.

\paragraph{Inference-aware scaling.} Chinchilla minimises \emph{training} compute. If a model
will serve $D_{\text{inf}}$ tokens, total cost is about $6ND + 2ND_{\text{inf}}$, and the optimum
moves to smaller models trained on more data (Sardana and Frankle, 2023). This is why open models
are routinely ``over-trained'' far past 20 tokens per parameter (\cref{sec:pretrain:scaling}).

\paragraph{Downstream metrics.} Loss scales smoothly; task accuracy often does not. Some
``emergent'' jumps are artefacts of discontinuous metrics such as exact match: a smoothly
improving per-token probability crosses a threshold (Schaeffer et al., 2023). Others are
genuinely hard to predict. Practical studies predict loss with a scaling law and map loss to
downstream metrics with a separate, empirical fit (Llama~3 reports such a two-step prediction).

\paragraph{Mixture-of-experts.} With MoE (\cref{sec:modern:moe}) total and active parameters
differ. Studies fit laws in active parameters and treat sparsity (experts per token over total
experts) as an extra axis. Several recent reports size their MoE models with IsoFLOP sweeps
over active parameters.

% ======================================================================
\section{Designing a budgeted study}\label{sec:scalinglaws:design}

Now the systems problem. You have a remote training service that accepts a model and recipe,
runs it on one accelerator, and returns validation losses. It charges a wall-clock budget
(say 12 device-hours) for fitting, and you must name the configuration and predict the loss for
a final run of 48 device-hours. This is the setting of any team that is allotted a slice of
a cluster to decide what to train next, and of Project~3 (\cref{ch:proj3}).

\subsection{The service model}\label{sec:scalinglaws:service}

\code{llmsys.scaling.SyntheticTrainingAPI} implements a service of this kind in-process. Its
semantics follow common shared-cluster policies:
\begin{itemize}
\item \code{submit(cfg, max\_runtime\_seconds)} validates the config, rejects exact duplicates,
      and \emph{reserves} the full \code{max\_runtime\_seconds} against your budget while the run
      is queued or running.
\item Runs execute FIFO on $W$ parallel workers in simulated time. When one finishes, the
      charge becomes its actual runtime, clipped to $[1, \text{max\_runtime}]$, and the
      unused reservation is refunded.
\item A run that needs longer than its reservation \emph{times out}: you are charged the full
      reservation and get back only the losses evaluated before the cut-off.
\item An LR far above the optimum diverges: losses end in \code{NaN}, and you pay for the
      time until divergence is detected.
\end{itemize}
Behind it sits a hidden loss surface: \cref{eq:scalinglaws:chin}, with $N$ the non-embedding
parameters, data discounted by large batches, a quadratic penalty in $\log(\text{LR}/\text{lr}^*(N))$,
a small aspect-ratio penalty, deterministic per-seed noise, and the mid-run cosine bump. Its
runtime model is the one in the Systems Lens of \cref{sec:scalinglaws:accounting}, with MFU
rising with width and batch.

\begin{figure}[t]
\centering
\begin{tikzpicture}[x=20cm, y=0.36cm, font=\scriptsize]
\fill[sysblue!25] (0.0000,-3.15) rectangle (0.0366,-3.85);
\fill[sysblue!80] (0.0000,-3.15) rectangle (0.0293,-3.85);
\node[anchor=east] at (0,-3.50) {1e+16};
\fill[sysblue!25] (0.0000,-4.15) rectangle (0.0337,-4.85);
\fill[sysblue!80] (0.0000,-4.15) rectangle (0.0270,-4.85);
\node[anchor=east] at (0,-4.50) {1e+16};
\fill[sysblue!25] (0.0000,-5.15) rectangle (0.0304,-5.85);
\fill[sysblue!80] (0.0000,-5.15) rectangle (0.0243,-5.85);
\node[anchor=east] at (0,-5.50) {1e+16};
\fill[sysblue!25] (0.0243,-6.15) rectangle (0.0537,-6.85);
\fill[sysblue!80] (0.0243,-6.15) rectangle (0.0478,-6.85);
\node[anchor=east] at (0,-6.50) {1e+16};
\fill[sysblue!25] (0.0270,-7.15) rectangle (0.1243,-7.85);
\fill[sysblue!80] (0.0270,-7.15) rectangle (0.1049,-7.85);
\node[anchor=east] at (0,-7.50) {4e+16};
\fill[sysblue!25] (0.0293,-8.15) rectangle (0.1134,-8.85);
\fill[sysblue!80] (0.0293,-8.15) rectangle (0.0966,-8.85);
\node[anchor=east] at (0,-8.50) {4e+16};
\fill[sysblue!25] (0.0478,-9.15) rectangle (0.1278,-9.85);
\fill[sysblue!80] (0.0478,-9.15) rectangle (0.1118,-9.85);
\node[anchor=east] at (0,-9.50) {4e+16};
\fill[sysblue!25] (0.0966,-10.15) rectangle (0.1708,-10.85);
\fill[sysblue!80] (0.0966,-10.15) rectangle (0.1560,-10.85);
\node[anchor=east] at (0,-10.50) {4e+16};
\fill[sysblue!25] (0.1049,-11.15) rectangle (0.4040,-11.85);
\fill[sysblue!80] (0.1049,-11.15) rectangle (0.3441,-11.85);
\node[anchor=east] at (0,-11.50) {2e+17};
\fill[sysblue!25] (0.1118,-12.15) rectangle (0.3944,-12.85);
\fill[sysblue!80] (0.1118,-12.15) rectangle (0.3378,-12.85);
\node[anchor=east] at (0,-12.50) {2e+17};
\fill[sysblue!25] (0.1560,-13.15) rectangle (0.4153,-13.85);
\fill[sysblue!80] (0.1560,-13.15) rectangle (0.3635,-13.85);
\node[anchor=east] at (0,-13.50) {2e+17};
\fill[sysblue!25] (0.3378,-14.15) rectangle (0.5759,-14.85);
\fill[sysblue!80] (0.3378,-14.15) rectangle (0.5283,-14.85);
\node[anchor=east] at (0,-14.50) {2e+17};

\draw[->] (0,0.2) -- (0.56,0.2) node[above left] {simulated hours};
\foreach \t in {0,0.1,0.2,0.3,0.4,0.5} \draw (\t,0.2) -- (\t,0) node[below] {\t};
\fill[sysblue!25] (0.36,-2.2) rectangle (0.38,-1.6); \node[anchor=west] at (0.385,-1.9) {reserved, refunded};
\fill[sysblue!80] (0.36,-3.2) rectangle (0.38,-2.6); \node[anchor=west] at (0.385,-2.9) {actual runtime};
\end{tikzpicture}
\caption{A 12-run IsoFLOP grid ($10^{16}$, $4\times10^{16}$, $1.6\times10^{17}$ FLOPs, four sizes
each) executed by the synthetic service on three workers. Light bars are each run's reservation
($1.25\times$ the predicted runtime), dark bars its actual runtime; the difference is refunded
at completion, but while the run is in flight it is unavailable for new submissions. Data:
simulated by \code{llmsys.scaling}.}
\label{fig:scalinglaws:gantt}
\end{figure}

\begin{systemslens}[Reservations are a concurrency limit]
Because queued and running jobs hold their full reservation, the number of experiments you can
have \emph{in flight} is bounded by $\text{remaining budget}/\text{reservation}$, not by the cluster.
Over-reserving (a safety factor of $3\times$ ``just in case'') does not cost budget in the end,
but it serialises your study and stretches it over more calendar time. Under-reserving costs
budget: a timed-out run is charged in full and returns nothing usable. Predict runtimes from
your accounting and a measured MFU, add a modest margin ($1.2$--$1.5\times$), and calibrate the
margin on the first few completed runs.
\end{systemslens}

\subsection{A staged plan}\label{sec:scalinglaws:stages}

A study that spends its budget well proceeds in stages, each small enough to revise the next:
\begin{enumerate}
\item \textbf{Dry run and accounting.} Submit the smallest sensible config once. Check that the
      service accepts your configs, that the losses look sane ($\ln V$ at initialisation, a
      steady fall), and that the measured runtime matches your prediction. This costs minutes
      and catches unit errors that would waste hours.
\item \textbf{Hyperparameter transfer.} Sweep the LR (and, if you can afford it, the batch) at
      three or four small sizes; fit $\text{lr}^*(N)$. Fix the aspect ratio by a rule.
\item \textbf{IsoFLOP grid.} Choose 3--5 budgets spaced by $3$--$4\times$, 5--7 sizes each,
      centred on your current best guess, wide enough ($8$--$16\times$ in $N$) to bracket the
      minimum. Use the fitted LR for each size.
\item \textbf{Fit and diagnose.} Fit methods 2 and 3; check residuals; bootstrap the allocation
      exponent. If a budget's minimum is at the edge of its grid, spend more there.
\item \textbf{Verification.} Spend part of the remaining budget on one run at the largest budget
      you can afford, at the predicted optimum. Compare its loss with the prediction. This is the
      single most informative run in the study.
\item \textbf{Final choice.} Choose the final config by searching configurations that
      \emph{finish} in the target wall-clock time (MFU varies with shape) under the fitted law,
      and report the predicted loss with an interval.
\end{enumerate}

\begin{napkinmath}[Where a geometric grid spends its budget]
With $m$ sizes per budget and budgets $C_{\max}, C_{\max}/r, C_{\max}/r^2,\dots$, the FLOPs spent
are $mC_{\max}(1+1/r+1/r^2+\dots)\approx mC_{\max}\,r/(r-1)$. With $r=4$ the largest budget is
$75\%$ of the grid, and with $r=3$ it is $67\%$. The top budget is what you are really buying,
so choose it first: given 12 hours, a 7-run top budget at $\sim$55\% of the study leaves room
for the smaller budgets, the LR sweep and a verification run. Then check the leverage: your
largest budget sits $\log_{10}(C_{\text{final}}/C_{\max})$ decades below the target, and every
decade of extrapolation multiplies exponent error into $N$ error.
\end{napkinmath}

\begin{lstlisting}[caption={Planning an IsoFLOP grid and pricing it in wall-clock seconds (from \code{llmsys/scaling.py}).},label={lst:scalinglaws:plan}]
def plan_isoflop_grid(budgets, sizes_per_budget=5, span=8.0,
                      hardware=SyntheticHardware(), seconds_available=math.inf,
                      safety=1.25, make_config=config_for_params, n_evals=8):
    plan, spent = [], 0.0
    for c in sorted(budgets):                 # cheapest budgets first
        n_mid = math.sqrt(c / (6 * 20))       # centre on 20 tokens/param
        ns = np.geomspace(n_mid / math.sqrt(span), n_mid * math.sqrt(span),
                          sizes_per_budget)
        rows = []
        for n in ns:
            cfg = make_config(n)              # shape + fitted LR
            cfg = replace(cfg, tokens=tokens_for_budget(cfg, c))
            rows.append((cfg, safety * hardware.runtime(cfg, n_evals)))
        cost = sum(r for _, r in rows)
        if spent + cost > seconds_available:
            break                             # never plan past budget
        plan.extend(rows); spent += cost
    return plan
\end{lstlisting}

\subsection{Choosing the final run when the budget is wall-clock}\label{sec:scalinglaws:wallchoice}

The closed form (\cref{eq:scalinglaws:closed}) answers ``best $N$ for $C$ FLOPs''. A
wall-clock budget answers a different question, because $C$ itself depends on the shape through
MFU. The robust approach is a search. Enumerate candidate sizes and batch sizes, compute how
many tokens each can process in the time limit at its MFU, predict its loss with the fitted
surface, and keep the best that finishes (\cref{lst:scalinglaws:choose}).
\Cref{fig:scalinglaws:wall} shows the difference on the synthetic hardware.

\begin{lstlisting}[caption={Wall-clock-aware final choice (Lab 12, Task 9).},label={lst:scalinglaws:choose}]
def choose_final_config(fit, lr_law, seconds, hw, sizes=None,
                        batch_sizes=(128, 256, 512), aspect=64.0):
    sizes = sizes if sizes is not None else np.geomspace(1e7, 3e9, 60)
    k, a = lr_law
    best = None
    for n in sizes:
        for bs in batch_sizes:
            cfg = config_for_params(n, batch_size=bs, aspect=aspect)
            n_ne = cfg.non_embedding_params()
            cfg = replace(cfg, lr=k * n_ne ** a)
            flops = hw.flops_in(cfg, seconds)   # depends on MFU(shape)
            cfg = replace(cfg, tokens=tokens_for_budget(cfg, flops))
            if hw.runtime(cfg) > seconds:
                continue                    # must actually finish
            pred = float(fit.loss(n_ne, cfg.tokens))
            if best is None or pred < best[1]:
                best = (cfg, pred)
    return best
\end{lstlisting}

\begin{figure}[t]
\centering
\begin{tikzpicture}
\begin{groupplot}[group style={group size=2 by 1, horizontal sep=1.7cm},
  width=0.49\linewidth, height=5cm, xmode=log, label style={font=\small},
  tick label style={font=\scriptsize}, xlabel={non-embedding params $N$}]
\nextgroupplot[ymode=log, ylabel={training FLOPs in 48 h}, ymin=1e19, ymax=1e20,
  legend style={font=\scriptsize, draw=none, at={(0.5,1.02)}, anchor=south, legend columns=2}]
\addplot[very thick, sysblue] coordinates {(8.847e+05,2.304e+19) (1.327e+06,2.304e+19) (1.327e+06,2.304e+19) (2.408e+06,2.764e+19) (2.408e+06,2.764e+19) (2.408e+06,2.764e+19) (3.211e+06,2.764e+19) (4.833e+06,3.141e+19) (4.833e+06,3.141e+19) (7.078e+06,3.455e+19) (8.847e+06,3.455e+19) (8.847e+06,3.455e+19) (1.219e+07,3.721e+19) (1.462e+07,3.721e+19) (1.868e+07,3.949e+19) (2.179e+07,3.949e+19) (2.787e+07,4.146e+19) (3.185e+07,4.146e+19) (3.965e+07,4.319e+19) (5.312e+07,4.472e+19) (6.37e+07,4.607e+19) (7.078e+07,4.607e+19) (9.196e+07,4.728e+19) (1.053e+08,4.837e+19) (1.327e+08,4.936e+19) (1.644e+08,5.026e+19) (1.979e+08,5.108e+19) (2.673e+08,5.252e+19) (3.133e+08,5.316e+19) (3.685e+08,5.375e+19) (4.663e+08,5.481e+19) (5.662e+08,5.528e+19) (6.955e+08,5.615e+19) (8.889e+08,5.691e+19) (1.062e+09,5.759e+19) (1.305e+09,5.819e+19) (1.595e+09,5.874e+19) (2.023e+09,5.946e+19) (2.392e+09,5.989e+19) (2.948e+09,6.047e+19)};
\addplot[thick, napkin, dashed, domain=1e6:3e9, samples=2] {7.774e+19};
\legend{achievable, at saturated MFU}
\nextgroupplot[ylabel={MFU}, ymin=0, ymax=0.45]
\addplot[very thick, labpurple] coordinates {(8.847e+05,0.1333) (1.327e+06,0.1333) (1.327e+06,0.1333) (2.408e+06,0.16) (2.408e+06,0.16) (2.408e+06,0.16) (3.211e+06,0.16) (4.833e+06,0.1818) (4.833e+06,0.1818) (7.078e+06,0.2) (8.847e+06,0.2) (8.847e+06,0.2) (1.219e+07,0.2154) (1.462e+07,0.2154) (1.868e+07,0.2286) (2.179e+07,0.2286) (2.787e+07,0.24) (3.185e+07,0.24) (3.965e+07,0.25) (5.312e+07,0.2588) (6.37e+07,0.2667) (7.078e+07,0.2667) (9.196e+07,0.2737) (1.053e+08,0.28) (1.327e+08,0.2857) (1.644e+08,0.2909) (1.979e+08,0.2957) (2.673e+08,0.304) (3.133e+08,0.3077) (3.685e+08,0.3111) (4.663e+08,0.3172) (5.662e+08,0.32) (6.955e+08,0.325) (8.889e+08,0.3294) (1.062e+09,0.3333) (1.305e+09,0.3368) (1.595e+09,0.34) (2.023e+09,0.3442) (2.392e+09,0.3467) (2.948e+09,0.35)};
\end{groupplot}
\end{tikzpicture}
\caption{The synthetic hardware model. Left: training FLOPs each model size can actually spend
in a 48-hour slot (blue) versus what a constant, saturated MFU would give (orange, dashed).
Right: the MFU curve behind it. Narrow models lose up to about half their FLOPs to low MFU, so a
wall-clock budget penalises them relative to a FLOP budget. Near this study's optimum
(about $5\times10^8$ parameters) the curve has flattened, and the wall-clock and FLOP optima almost
coincide; for smaller budgets they do not. Simulated, not measured.}
\label{fig:scalinglaws:wall}
\end{figure}

\subsection{How sure are we? Bootstrapping the extrapolation}\label{sec:scalinglaws:uncertainty}

Every number in a scaling study comes from a few dozen noisy runs, and the quantity you care
about (the size at $10^{20}$ FLOPs) is an extrapolation. Quantify the uncertainty by
resampling. Resample runs with replacement \emph{within each budget}, which keeps the IsoFLOP
structure; redo the whole fit; and collect the predictions (\cref{fig:scalinglaws:fan}). The
spread grows with distance from the data, at a rate set by the exponent's uncertainty.

\begin{figure}[t]
\centering
\begin{tikzpicture}
\begin{loglogaxis}[width=0.72\linewidth, height=5.4cm, xmin=1e16, xmax=1e20,
  xlabel={compute $C$ (FLOPs)}, ylabel={$N_{\text{opt}}$}, label style={font=\small},
  tick label style={font=\scriptsize},
  legend style={font=\scriptsize, draw=none, at={(0.02,0.98)}, anchor=north west}]
\addplot[fill=sysblue!20, draw=none, forget plot] coordinates {(1e+16,1.425e+07) (2.154e+16,1.914e+07) (4.642e+16,2.594e+07) (1e+17,3.585e+07) (2.154e+17,5.065e+07) (4.642e+17,7.676e+07) (1e+18,1.124e+08) (2.154e+18,1.647e+08) (4.642e+18,2.419e+08) (1e+19,3.556e+08) (2.154e+19,5.227e+08) (4.642e+19,7.684e+08) (1e+20,1.13e+09) (1e+20,4.109e+08) (4.642e+19,3.125e+08) (2.154e+19,2.372e+08) (1e+19,1.794e+08) (4.642e+18,1.355e+08) (2.154e+18,1.019e+08) (1e+18,7.702e+07) (4.642e+17,5.756e+07) (2.154e+17,4.173e+07) (1e+17,3.044e+07) (4.642e+16,2.031e+07) (2.154e+16,1.33e+07) (1e+16,9.533e+06)} -- cycle;
\addplot[very thick, sysblue] coordinates {(1e+16,1.297e+07) (2.154e+16,1.796e+07) (4.642e+16,2.475e+07) (1e+17,3.391e+07) (2.154e+17,4.68e+07) (4.642e+17,6.444e+07) (1e+18,8.9e+07) (2.154e+18,1.228e+08) (4.642e+18,1.695e+08) (1e+19,2.34e+08) (2.154e+19,3.229e+08) (4.642e+19,4.447e+08) (1e+20,6.131e+08)};
\addplot[thick, black, dashed] coordinates {(1e+16,9.577e+06) (2.154e+16,1.356e+07) (4.642e+16,1.92e+07) (1e+17,2.719e+07) (2.154e+17,3.849e+07) (4.642e+17,5.45e+07) (1e+18,7.717e+07) (2.154e+18,1.093e+08) (4.642e+18,1.547e+08) (1e+19,2.191e+08) (2.154e+19,3.102e+08) (4.642e+19,4.392e+08) (1e+20,6.219e+08)};
\addplot[only marks, mark=*, warnred] coordinates {(9.996e+15,1.313e+07) (2.999e+16,2.08e+07) (9.999e+16,3.387e+07) (3e+17,5.385e+07)};
\legend{bootstrap median, hidden truth, IsoFLOP minima}
\end{loglogaxis}
\end{tikzpicture}
\caption{A 12-hour synthetic study (four budgets up to $3\times10^{17}$, six sizes each) and a
150-sample bootstrap. The band is the 5--95\% interval of the extrapolated $N_{\text{opt}}$;
at $10^{20}$ FLOPs it spans $4.1\times10^8$--$11.3\times10^8$ parameters, almost a factor of three, from
a fit whose individual points look very clean. Simulated with \code{llmsys.scaling}.}
\label{fig:scalinglaws:fan}
\end{figure}

\begin{napkinmath}[Exponent error becomes size error]
If the fitted exponent is off by $\delta a$, the predicted size at a budget $k$ decades beyond
your data is off by a factor $10^{k\,\delta a}$. With $\delta a=0.03$ (a typical bootstrap
standard deviation for a small study) and $k=2.5$ decades, that is $10^{0.075}\approx1.19$,
roughly a 20\% error in $N$. The good news is that IsoFLOP valleys are flat near the bottom: a model
20\% off the optimum loses only a small fraction of a percent of loss. The bad news is that the
\emph{loss prediction} inherits the error in $E$, which is the least well-determined parameter of
method 3. Predict the loss from the fit, then sanity-check it against the trend of your IsoFLOP minima
$L_{\text{opt}}(C)$.
\end{napkinmath}

\subsection{Reading a study log}\label{sec:scalinglaws:log}

\Cref{lst:scalinglaws:log} is the output of the reference solution of Lab~12 on the default
synthetic service: 12 simulated hours of fitting budget, a 48-hour target. Read it the way you
would read a real study.

\begin{lstlisting}[style=plain, caption={Output of \code{run.py study} (synthetic service, seed 0; numbers abridged).},label={lst:scalinglaws:log}]
lr law      : lr*(N) = 0.341 * N^-0.255
fit         : E=1.769 A=513.8 B=493.2 alpha=0.356 beta=0.290
allocation  : N_opt ~ C^0.407   D_opt ~ C^0.749   L_opt ~ C^-0.116
final config: L=23 d=1472 heads=23 d_ff=3904 batch=512 lr=1.98e-3
              tokens=1.46e10  (N_ne = 0.60e9, ~24 tokens/param)
predicted   : 2.7509   (after verification residual)
budget      : used 32,925 s of 43,200 s
score       : true 2.7453   oracle 2.7412   regret 0.0041   pred. error +0.0056
\end{lstlisting}

Four things to notice. The learning-rate law falls with size (exponent $-0.25$; the hidden truth
uses $-0.30$), so a single LR for all sizes would have been wrong at both ends. The parametric fit
recovers the hidden exponents to within $0.01$ ($0.35$ and $0.29$ in the truth), while the
IsoFLOP allocation exponents do not add to one ($0.41+0.75$). That gap is the accounting effect
of \cref{sec:scalinglaws:accounting}: $D$ was computed with the full FLOP count, including the LM
head and attention, while $N$ is non-embedding, so the two power laws are fits of different
identities. The choice itself is excellent (regret $0.004$ nats). Finally, before the
verification run the same fit predicted $2.711$, an error of $-0.034$. One run at the largest
affordable scale cut that error by a factor of six, because $E$ was the least constrained
parameter of the fit and the verification residual corrects it directly. A study without a
verification run reports a loss it has never seen.

\begin{pitfall}[Five ways scaling studies go wrong]
\begin{enumerate}
\item The grid does not bracket the minimum at the largest budget, which is the one that matters most.
\item One learning rate is used for every size.
\item Intermediate (mid-schedule) losses are fitted as if they were finished runs.
\item $ and $ use inconsistent definitions across runs or between the fit and the final config.
\item The final prediction is reported with no interval and no verification.
\end{enumerate}
\end{pitfall}

% ======================================================================
\section{Running real sweeps}\label{sec:scalinglaws:running}

The synthetic service lets you practise the analysis. On real hardware the same study is mostly
an exercise in orchestration. A layout that scales from a laptop to a cluster:

\begin{lstlisting}[style=plain]
study/
  configs/          one YAML/JSON per run, generated by a script (never hand-edited)
  ledger.jsonl      append-only: submitted request, job id, result, timestamps
  runs/<run_id>/    checkpoints, metrics.jsonl, stdout, the exact config + git hash
  analysis/         fits, bootstrap draws, plots -- regenerated from ledger.jsonl
  plan.md           the staged plan and every decision, with dates
\end{lstlisting}

Three rules make this robust. \textbf{Generate configs from code.} The planner of
\cref{lst:scalinglaws:plan} writes the files, and its output is versioned with the study.
\textbf{Make submission idempotent.} Key every run by a hash of its canonical config, and refuse to
resubmit a key that is already in the ledger. A crashed notebook or a re-run cell then costs
nothing (the \code{Ledger} in \code{code/projects/p3\_scaling/ledger.py} does exactly this).
\textbf{Derive all analysis from the ledger}, never from a notebook's in-memory state.

On a Slurm cluster, an IsoFLOP grid is a natural job array: one array per budget, one task per
size, each task reading its config by index. This gives every task the same time limit, which
corresponds to the reservation above.

\begin{lstlisting}[style=shell, caption={One IsoFLOP budget as a Slurm job array (each task trains one size).},label={lst:scalinglaws:slurm}]
#!/bin/bash
#SBATCH --job-name=iso-1e18
#SBATCH --array=0-6                 # seven sizes
#SBATCH --gres=gpu:1 --cpus-per-task=8
#SBATCH --time=02:30:00             # reservation = 1.25 x predicted runtime
#SBATCH --output=runs/%x-%a.out
CFG=$(ls configs/iso-1e18/*.json | sed -n "$((SLURM_ARRAY_TASK_ID+1))p")
srun python train.py --config "$CFG" --metrics runs/$(basename "$CFG" .json)/metrics.jsonl
\end{lstlisting}

Log every run to one experiment tracker project (for example W\&B, or the JSONL logger of
\cref{ch:optim}), grouped by budget, with the config hash as the run name. Then the
IsoFLOP plot is a query, not a manual export. Seed the model initialisation but keep the data
order fixed across runs, so that differences between sizes are not differences between data
samples. Finally, record \emph{achieved throughput} for every run: it is the input to the
wall-clock planning of the next study. Lab 12's \code{run.py cpu-sweep} is a laptop-sized version
of this pipeline that trains real byte-level Qwen3-shaped models and fits what it measures.

% ======================================================================
\section{Summary}\label{sec:scalinglaws:summary}
\begin{itemize}
\item Scaling laws are power laws plus an irreducible term; they are measured, not derived, and
      their constants depend on data, architecture, optimizer and accounting choices.
\item Account consistently: non-embedding $N$, $C/D\approx6(N_{\text{ne}}+Vd)+6LTd$, and actual
      (not nominal) compute per run.
\item Kaplan and Chinchilla disagree mostly because of schedules, accounting and scale; the
      allocation exponent is sensitive, and its errors grow with extrapolation distance.
\item Fit allocation laws with IsoFLOP minima (bracketed, min or vertex) and with the parametric
      surface (Huber in log space, many starts); compare them, and bootstrap.
\item Model or remove scale dependence in LR and batch size: fit $\text{lr}^*(N)$ and respect
      the critical batch, or use $\mu$P.
\item WSD schedules with branched decays make durations cheap, if the checkpointing
      infrastructure is right.
\item A study is a budgeted experiment on a cluster: reservations limit concurrency, MFU makes
      wall-clock differ from FLOPs, and ledgers, staged plans and a verification run keep it
      honest.
\end{itemize}

% ======================================================================
\section{Exercises}\label{sec:scalinglaws:exercises}

\begin{exercise}[Closed-form allocation]\label{ex:scalinglaws:lagrange}
Derive \cref{eq:scalinglaws:closed} by minimising \cref{eq:scalinglaws:chin} subject to
$6ND=C$. Show that at the optimum the two reducible terms are in the fixed ratio
$\alpha A N^{-\alpha} : \beta B D^{-\beta} = 1:1$, and interpret that ratio.
\end{exercise}

\begin{exercise}[IsoFLOP minima lie on a power law]
Show that the minima of the IsoFLOP profiles of \cref{eq:scalinglaws:chin} satisfy
$N_{\text{opt}}\propto C^{\beta/(\alpha+\beta)}$ exactly. Then add an embedding term to the FLOP
count, $C=6(N+Vd)D$, and argue qualitatively how the measured exponent changes at small scale.
\end{exercise}

\begin{exercise}[Accounting for your own models]
For Qwen3-0.6B (\cref{ch:modern}) compute $N_{\text{ne}}$, the embedding fraction, and
$C/D$ from \cref{eq:scalinglaws:fpt} at $T=4096$. How large is the attention term relative to
$6N$? Which of the 12$Ld^2$ assumptions does Qwen3 violate (GQA, $\dff/d$, head dimension)?
\end{exercise}

\begin{exercise}[Quantisation error of the min method]
Sizes are spaced by a factor $r$ on a valley that is locally parabolic in $\log N$ with
curvature $\kappa$. Bound the error in $\log N_{\text{opt}}$ and in $L_{\text{opt}}$ of the
min-of-runs reading, and compare it with the noise level at which a parabola fit stops helping.
\end{exercise}

\begin{exercise}[Critical batch arithmetic]
From \cref{eq:scalinglaws:critbatch}, show that training with batch $B$ needs
$S=S_{\min}(1+B_{\text{crit}}/B)$ steps and $E=E_{\min}(1+B/B_{\text{crit}})$ examples. If a
step takes time $t_0 + B/\text{throughput}$, find the batch that minimises wall-clock time to the
target loss, and explain why it is larger than $B_{\text{crit}}$ when $t_0$ is large.
\end{exercise}

\begin{exercise}[Mid-run losses are not finished runs]\label{ex:scalinglaws:midrun}
Using \code{SyntheticTrainingAPI}, train one size for a long run with 16 evaluations and treat the
intermediate losses as runs of shorter length. Fit the data exponent $\beta$ from them and
compare it with a fit from separate runs of each length. Explain the bias using the cosine
schedule, and say how WSD branching would remove it.
\end{exercise}

\begin{exercise}[Cost of a WSD study]
For $m$ sizes and $k$ log-spaced durations (ratio 2), compare the total tokens of a cosine study
(one run per pair) with a WSD study (one longest run per size plus decays of a fraction $f$ of
each length). Give the asymptotic saving as $k\to\infty$.
\end{exercise}

\begin{exercise}[$\mu$P in one layer]
For a linear layer $y=Wx$ with $x\in\R^n$, $\|x\|_\infty=\Theta(1)$, show that an Adam update
(entries of magnitude $\approx\eta$) aligned with $x$ changes $y$ by $\Theta(\eta n)$. Derive the
$\eta\propto1/n$ rule, and explain why the same argument gives a different rule for SGD.
\end{exercise}

\begin{exercise}[Inference-aware optimum]
Minimise $\hat L(N,D)$ subject to $6ND + 2ND_{\text{inf}} = C_{\text{total}}$ for fixed
$D_{\text{inf}}$. Show that the optimal $N$ is smaller than the Chinchilla optimum, and compute
the optimal $N$ and tokens-per-parameter ratio for the default synthetic truth
(\code{GroundTruth()}) at $C_{\text{total}}=10^{21}$ and $D_{\text{inf}}=10^{11}$.
\end{exercise}

\begin{exercise}[Reservation policy]
Runtime predictions have a log-normal error with $\sigma=0.15$. A timed-out run is charged its
reservation $s\cdot\hat t$ and returns nothing; a completed run is charged its true runtime.
Find the safety factor $s$ that minimises the expected budget per useful run, and compare it
with the factor that minimises calendar time when at most $R/(s\hat t)$ runs can be in flight.
\end{exercise}

\begin{exercise}[Replicate the study]
Run Lab 12's \code{run.py replicate --seeds 10}. Report the
distribution of regret and prediction error. Which stage of your plan would you change to
reduce the spread, and what would it cost?
\end{exercise}

% ======================================================================
\section{Lab 12: A budgeted scaling study}\label{sec:scalinglaws:lab}

\begin{labbox}{12}{A budgeted scaling study against a synthetic training service}
\labdir{lab12_scaling_laws}
\labsection{Goals}
Build every piece of a scaling study (accounting, power-law and parametric fits, LR transfer,
uncertainty, wall-clock-aware choice) and run it end to end against a service with a hidden
ground truth, a wall-clock budget, reservations, timeouts and divergence.

\labsection{Setup}
\begin{lstlisting}[style=shell]
cd code
LAB_IMPL=starter  python -m pytest labs/lab12_scaling_laws    # your work
LAB_IMPL=solution python -m pytest labs/lab12_scaling_laws    # reference
\end{lstlisting}
The service is \code{llmsys.scaling.SyntheticTrainingAPI}; it is provided, not a task.

\labsection{Tasks}
\begin{enumerate}
\item \ident{count_params_and_flops}: \cref{eq:scalinglaws:12ld2,eq:scalinglaws:fpt}; the
      test checks your count against an actual \code{Qwen3Model}.
\item \ident{fit_power_law}: least squares in log-log space.
\item \ident{isoflop_optima}: min-of-runs and bracketed parabola vertex per budget.
\item \ident{fit_allocation}: $N_{\text{opt}}(C)$, $D_{\text{opt}}(C)$, $L_{\text{opt}}(C)$.
\item \ident{fit_chinchilla}: Huber/LSE fit from a grid of starts.
\item \ident{chinchilla_optimal}: \cref{eq:scalinglaws:closed}.
\item \ident{fit_lr_optimum}, \ident{fit_lr_scaling}: LR basin vertex and $\text{lr}^*(N)$.
\item \ident{bootstrap_ci}: percentile intervals for any fitted quantity.
\item \ident{choose_final_config}: best config that finishes in the time limit.
\item \ident{run_study}: the staged plan of \cref{sec:scalinglaws:stages}, within budget,
      including a verification run whose residual corrects the loss prediction, ending with
      \code{api.submit\_final}. The test requires low regret against the oracle.
\end{enumerate}

\labsection{Measurements}
\begin{itemize}
\item \code{run.py study --plot study.png}: your IsoFLOP profiles and allocation fit; compare
      with \cref{fig:scalinglaws:isoflop}.
\item \code{run.py replicate --seeds 10}: regret and prediction-error distributions
      (\cref{ex:scalinglaws:midrun} and the replication exercise).
\item Ablate: a fixed LR for all sizes instead of $\text{lr}^*(N)$; min versus vertex readings;
      no verification run. Which ablation hurts regret most? Which hurts the loss prediction?
\item \code{run.py cpu-sweep}: a real IsoFLOP sweep with tiny models. Expect noisy exponents and
      explain where the noise comes from (data, seeds, the overhead-bound MFU of tiny models).
\end{itemize}

\labsection{Deliverables}
Passing tests; a one-page write-up with your plan, budget usage, both fits, a bootstrap
interval for the final $N$ and loss, and the score from \code{api.score}.

\labsection{Stretch goals}
Add WSD branching to the synthetic service (a \code{branch(eid, at\_fraction)} endpoint that
charges only the decay) and redesign the study around method~1; add a batch-size sweep and fit
$B_{\text{opt}}(C)$.
\end{labbox}
